\chapter{M004 --- Surface Termination Selection}

\section{Stage 1 --- Mathematical Specification}

\subsection{Scientific Purpose}

Select a single surface termination from an admissible Termination Set.

\subsection{Transformation}

\[
\Phi_{\mathrm{select}}:
\mathcal{P}(\mathcal{T})
\rightarrow
\mathcal{T}
\]

\subsection{Mathematical Inputs}

\begin{itemize}
\item Termination Set
\item Selection criterion
\end{itemize}

\subsection{Mathematical Output}

A single Surface Termination.

\subsection{Domain}

Finite sets of admissible surface terminations.

\subsection{Codomain}

Surface-termination space.

\subsection{Preconditions}

\begin{itemize}
\item Non-empty Termination Set.
\item Well-defined selection criterion.
\end{itemize}

\subsection{Postconditions}

A selected termination together with its provenance.

\subsection{Invariants}

\begin{itemize}
\item Primitive surface identity.
\item Crystallographic orientation.
\item Termination provenance.
\end{itemize}

\subsection{Gauge Freedoms}

None in the selection itself. Gauge freedoms have already been quotiented
out during termination enumeration.

\subsection{Selection Criteria}

Selection may be driven by reproducibility, user choice, symmetry,
stoichiometry, polarity, or other explicitly specified admissibility
criteria. Such criteria form part of the mathematical contract of this
transformation.
